% =====================================================================
%  THÈME 4 — Réactif limitant et avancement
%  Niveau MOYEN — Corrigé
% =====================================================================
\documentclass[11pt,a4paper]{article}
\usepackage[utf8]{inputenc}
\usepackage[T1]{fontenc}
\usepackage{lmodern}
\usepackage[french]{babel}
\usepackage{amsmath,amssymb,mathtools}
\usepackage{icomma}
\usepackage[version=4]{mhchem}
\usepackage{siunitx}
\sisetup{output-decimal-marker={,},per-mode=symbol}
\usepackage{xcolor}
\usepackage{tikz}
\usepackage[most]{tcolorbox}
\usepackage{enumitem}
\usepackage{array}
\usepackage{booktabs}
\usepackage{fancyhdr}
\usepackage[a4paper,margin=2.1cm,headheight=26pt,headsep=8pt]{geometry}

\definecolor{primary}{RGB}{142,93,168}
\definecolor{primarydark}{RGB}{82,45,108}
\definecolor{corrcol}{RGB}{176,110,20}
\newcommand{\nq}[1]{n_{\text{#1}}}
\newcommand{\Vm}{V_{\mathrm m}}

\pagestyle{fancy}\fancyhf{}
\fancyhead[L]{\small\textcolor{primarydark}{\textbf{Fiche 5} — Thème 4 : Limitant}}
\fancyhead[R]{\small\textcolor{corrcol}{\textsc{Corrigé — Moyen}}}
\fancyfoot[C]{\small\thepage}
\renewcommand{\headrulewidth}{0.8pt}
\renewcommand{\headrule}{\hbox to\headwidth{\color{primary}\leaders\hrule height \headrulewidth\hfill}}

\newtcolorbox[auto counter]{cor}[1][]{enhanced,breakable,boxrule=1pt,arc=3pt,
  left=3mm,right=3mm,top=2.4mm,bottom=2.4mm,colback=corrcol!6,colframe=corrcol,
  fonttitle=\bfseries,coltitle=white,
  title={Corrigé~\thetcbcounter\ \ifx\relax#1\relax\relax\else\textmd{--- \itshape #1}\fi}}

\newcommand{\rep}[1]{\par\smallskip\noindent\textbf{\color{corrcol}$\Rightarrow$ #1}}

\setlength{\parindent}{0pt}
\setlength{\parskip}{3pt}

\begin{document}

\begin{center}
{\LARGE\bfseries\color{primarydark}Thème 4 — Réactif limitant}\\[1pt]
{\large\color{corrcol}\textbf{Corrigé — Niveau Moyen}}\\
{\color{primary}\rule{0.6\linewidth}{1pt}}
\end{center}

\begin{cor}[bilan complet par le tableau d'avancement]
\textbf{a)} Tableau :
\begin{center}\renewcommand{\arraystretch}{1.25}\small
\begin{tabular}{l|cccc}
 & $2\,\ce{A}$ & $3\,\ce{B}$ & \ce{C} & $6\,\ce{D}$\\ \hline
$\xi=0$ & $0{,}100$ & $0{,}200$ & $0$ & $0$\\
$\xi$ & $0{,}100-2\xi$ & $0{,}200-3\xi$ & $\xi$ & $6\xi$\\
\end{tabular}
\end{center}
\textbf{b)} On annule chaque réactif : \ce{A} $\Rightarrow \xi=\frac{0{,}100}{2}=0{,}050$ ;
\ce{B} $\Rightarrow \xi=\frac{0{,}200}{3}\approx0{,}0667$. Le plus petit est $0{,}050$, donc
$\xi_{\max}=\SI{0,050}{\mole}$ et \boxed{\text{\ce{A} est limitant}}.\\
\textbf{c)} $\nq{C}=\xi_{\max}=\SI{0,050}{\mole}$ ; $\nq{D}=6\xi_{\max}=\SI{0,300}{\mole}$ ;
$\nq{B,\text{reste}}=0{,}200-3\times0{,}050=\boxed{\SI{0,050}{\mole}}$.
\end{cor}

\begin{cor}[trois réactifs]
\[
\frac{n(\ce{K2Cr2O7})}{2}=\frac{0{,}10}{2}=0{,}050 ;\quad
\frac{n(\ce{C})}{3}=\frac{0{,}20}{3}\approx0{,}0667 ;\quad
\frac{n(\ce{H2SO4})}{8}=\frac{0{,}30}{8}=0{,}0375 .
\]
Le plus petit est $0{,}0375$ : \boxed{\text{\ce{H2SO4} est le réactif limitant}}
($0{,}0375<0{,}050<0{,}0667$).
\end{cor}

\begin{cor}[zinc et acide : gaz dégagé et excès]
\textbf{a)} $n(\ce{Zn})=\dfrac{6{,}5}{65}=0{,}10$~mol ; $n(\ce{HCl})=1{,}5\times0{,}100=0{,}15$~mol.
Rapports : $\frac{0{,}10}{1}=0{,}10$ et $\frac{0{,}15}{2}=0{,}075$. Le plus petit est celui de \ce{HCl} :
\boxed{\text{\ce{HCl} limitant}}.\\
\textbf{b)} $\nq{H2}=n(\ce{HCl})\times\dfrac{1}{2}=0{,}15\times0{,}5=0{,}075$~mol ;
$V=n\Vm=0{,}075\times22{,}4=\boxed{\SI{1,68}{\litre}}$.\\
\textbf{c)} \ce{Zn} consommé $=n(\ce{HCl})\times\frac{1}{2}=0{,}075$~mol ;
reste $0{,}10-0{,}075=\boxed{\SI{0,025}{\mole}}$ de \ce{Zn}.
\end{cor}

\begin{cor}[combustion incomplète en réactifs]
\textbf{a)} $n(\ce{C3H8})=\dfrac{88}{44}=2{,}0$~mol ; $n(\ce{O2})=\dfrac{16}{32}=0{,}50$~mol.
Rapports : $\frac{2{,}0}{1}=2{,}0$ et $\frac{0{,}50}{5}=0{,}10$. \boxed{\text{\ce{O2} est limitant}}.\\
\textbf{b)} $\nq{CO2}=n(\ce{O2})\times\dfrac{3}{5}=0{,}50\times0{,}6=0{,}30$~mol ;
$V=0{,}30\times22{,}4=\boxed{\SI{6,72}{\litre}}$.
\rep{Malgré \SI{88}{\gram} de propane, le manque de \ce{O2} limite tout : le carburant est en large excès.}
\end{cor}

\begin{cor}[combustion du sulfure d'hydrogène]
\textbf{a)} $n(\ce{H2S})=\dfrac{10{,}2}{34}=0{,}30$~mol ; $n(\ce{O2})=\dfrac{9{,}6}{32}=0{,}30$~mol.
Rapports : $\frac{0{,}30}{2}=0{,}15$ et $\frac{0{,}30}{3}=0{,}10$. \boxed{\text{\ce{O2} est limitant}}.\\
\textbf{b)} $\nq{SO2}=n(\ce{O2})\times\dfrac{2}{3}=0{,}30\times\dfrac{2}{3}=0{,}20$~mol ;
$m=0{,}20\times64=\boxed{\SI{12,8}{\gram}}$.
\end{cor}

\end{document}
